\section{The Gravitational N-Body Problem}

The N-body problem models the dynamical evolution of $N$ point masses
interacting through Newtonian gravity. It is one of the oldest unsolved
problems in mathematical physics: for $N > 2$ there exists no general
closed-form solution, requiring numerical integration.

This simulation solves the equations of motion for thousands to millions
of gravitating particles, revealing the emergence of large-scale
structure: galaxy clusters, filaments, and voids.

\section{Gravitational Force Law}

Every pair of particles $(i, j)$ with masses $m_i$, $m_j$ exerts an
attractive force along the line connecting them. The softened force
law prevents singularities at close approach:

\begin{equation}
F_{ij} = -G \cdot m_i m_j / (|r_{ij}|^2 + \epsilon^2)
\end{equation}

\param{G}{Gravitational constant (coupling strength)}
\param{r_{ij}}{Displacement vector between particles $i$ and $j$}
\param{\epsilon}{Softening length (prevents close-encounter divergence)}

\section{Equations of Motion}

The Hamiltonian system has $6N$ degrees of freedom (3 positions and
3 velocities per particle). Newton's second law for particle $i$:

\begin{equation}
a_i = \sum_{j \neq i} G \cdot m_j \cdot r_{ij} / (|r_{ij}|^2 + \epsilon^2)^{3/2}
\end{equation}

The system is time-reversible and conserves total energy
$E = T + V$ (kinetic + potential) and total momentum $P = \sum m_i v_i$.

\section{Leapfrog Integration}

We use the symplectic leapfrog (Verlet) integrator. Unlike Runge-Kutta
methods, it exactly preserves the symplectic structure of Hamiltonian
mechanics, giving excellent long-term energy conservation.

The kick-drift-kick scheme:

\begin{equation}
v_{n+\frac{1}{2}} = v_n + a_n \cdot \Delta t / 2
\end{equation}

\begin{equation}
r_{n+1} = r_n + v_{n+\frac{1}{2}} \cdot \Delta t
\end{equation}

\begin{equation}
v_{n+1} = v_{n+\frac{1}{2}} + a_{n+1} \cdot \Delta t / 2
\end{equation}

This scheme is second-order accurate and time-reversible.

\section{Barnes-Hut Algorithm}

The brute-force approach computes all pairwise interactions with
$O(N^2)$ complexity. For large $N$ this becomes prohibitive.

The Barnes-Hut tree code groups distant particles into an octree and
approximates their collective gravity using multipole moments. A cell
is treated as a single body when:

\begin{equation}
\theta > s / d
\end{equation}

where $s$ is the cell size, $d$ is the distance to the particle, and
$\theta$ is the opening angle parameter. This reduces the complexity
to $O(N \log N)$, enabling simulations with $10^5$--$10^6$ particles
in real time on a GPU.

\section{Simulation Parameters}

\param{G}{Gravitational constant (coupling strength)}
\param{\Delta t}{Integration timestep (smaller = more accurate)}
\param{\epsilon}{Softening length (prevents close-encounter divergence)}
\param{\theta}{Barnes-Hut opening angle (accuracy vs. speed)}
\param{N}{Number of particles}

\section{References}

\begin{itemize}
\item Newton, I. (1687). Principia Mathematica.
\item Barnes, J. \& Hut, P. (1986). Nature 324, 446--449.
\item Springel, V. et al. (2005). Nature 435, 629--636.
\item Dehnen, W. \& Read, J. (2011). EPJP 126, 55.
\end{itemize}
